\section{Introdução}

Pipelines de dados em lote são reexecutados com frequência: o orquestrador repete uma tarefa que falhou, um engenheiro reprocessa um intervalo de datas (\emph{backfill}) ou dados chegam atrasados e exigem recálculo. Em todos esses casos, espera-se que a nova execução produza o mesmo resultado que uma execução única e correta produziria. Quando isso não ocorre, o defeito costuma ser silencioso: linhas duplicadas, totais inflados ou partições gravadas na data errada.

Na comunidade de engenharia de dados, a resposta mais difundida a esse problema é adotar um estilo funcional. \citet{beauchemin2018} propõe tratar cada tarefa como uma função pura que sobrescreve integralmente uma partição de saída, e a documentação do Apache Airflow recomenda que tarefas produzam o mesmo resultado em cada reexecução, evitem \texttt{INSERT} e não usem o relógio do sistema em cálculos críticos~\citep{airflow2023}. Essas recomendações, porém, agrupam práticas de naturezas diferentes. Algumas decorrem diretamente da transparência referencial, como receber a data como parâmetro e evitar efeitos colaterais. Outras dizem respeito à interação com o armazenamento ou à semântica dos dados, como sobrescrever partições e deduplicar registros, e podem ser adotadas em código imperativo sem nenhuma mudança de estilo. Na busca realizada para este trabalho, não foi encontrado estudo que separe experimentalmente a contribuição de cada grupo.

Este artigo investiga duas questões de pesquisa:
\begin{itemize}
  \item \textbf{QP1:} quais práticas, isoladamente e em combinação, determinam se um pipeline produz o resultado correto após reexecuções, \emph{backfills}, falhas parciais e chegadas irregulares de dados?
  \item \textbf{QP2:} as transformações puras tornam o pipeline mais testável?
\end{itemize}

As contribuições desta versão preliminar são: (i) um desenho fatorial que isola quatro práticas em um mesmo pipeline, com hipóteses registradas antes das medições; (ii) os resultados da QP1 para as 16 configurações em seis cenários; e (iii) um artefato reprodutível com o pipeline, o gerador de dados e o oráculo\footnote{Repositório: [LINK A DEFINIR].}. A QP2 está em andamento. Os PDFs dos trabalhos citados estão disponíveis em uma pasta compartilhada\footnote{Pasta com os PDFs citados: [LINK A DEFINIR].}.

\section{Fundamentação e Trabalhos Relacionados}

\paragraph{Pureza e composição.} Uma função é pura quando seu resultado depende apenas dos argumentos e sua avaliação não produz efeitos observáveis. \citet{hughes1989} argumenta que funções de ordem superior e avaliação preguiçosa são as características que mais contribuem para a modularidade de programas funcionais. A influência dessas ideias no processamento de dados em larga escala é registrada por \citet{hu2015}, que apontam o MapReduce como inspirado nas funções \texttt{map} e \texttt{reduce}.

\paragraph{Recomputação e determinismo.} Os \emph{Resilient Distributed Datasets} do Spark obtêm tolerância a falhas por meio de transformações de granularidade grossa, em vez de atualizações finas de estado compartilhado~\citep{zaharia2012}: uma partição perdida é recalculada reaplicando as transformações que a produziram. Essa recuperação só reproduz o valor original se as transformações forem determinísticas, o que transfere ao código do usuário a responsabilidade pela pureza.

\paragraph{Engenharia de dados funcional.} \citet{beauchemin2018} sistematiza a ideia de tarefas puras e partições imutáveis em pipelines em lote, e a documentação do Airflow traduz parte dela em regras operacionais: tratar tarefas como transações, nunca produzir resultados incompletos, ler e escrever em partições específicas e não usar \texttt{now()} em cálculos críticos~\citep{airflow2023}. Ambas as fontes são normativas: descrevem boas práticas, mas não medem seu efeito.

\paragraph{Núcleo funcional e testabilidade.} O padrão \emph{functional core, imperative shell}~\citep{bernhardt2012} concentra a lógica em funções puras e isola entrada, saída, banco de dados e rede numa casca imperativa; segundo o autor, as partes funcionais podem ser testadas facilmente e, muitas vezes, sem dublês de teste. O teste baseado em propriedades nasceu no contexto funcional~\citep{claessen2000} e está disponível em Python pela biblioteca Hypothesis~\citep{maciver2019}; o teste de mutação mede a eficácia de uma suíte pela proporção de defeitos injetados que ela detecta~\citep{jia2011}.

\paragraph{Qualidade de pipelines.} \citet{foidl2024} propõem uma taxonomia de fatores que influenciam a qualidade de pipelines de dados, validada com entrevistas e com a análise de 600 \emph{issues} de 11 projetos. O trabalho caracteriza problemas, mas não avalia práticas de programação específicas, que é o foco deste estudo.

\section{Desenho do Estudo}

O desenho completo, incluindo parâmetros dos dados e hipóteses, foi registrado no repositório antes de qualquer medição, em um \emph{commit} identificado por uma \emph{tag}.

\subsection{Objeto de Estudo}

O objeto é um pipeline em Python que processa os eventos de pedidos de uma data lógica \emph{D} e grava uma partição de saída para \emph{D}. As etapas são: leitura dos eventos de \emph{D} já disponíveis; normalização de dois formatos de registro (plano e aninhado), com descarte dos malformados; validação de status e valor; deduplicação; enriquecimento com o nível do cliente; agregação de pedidos e receita por nível; e escrita. As etapas de transformação formam um núcleo funcional, e relógio, origem e destino ficam numa casca de efeitos, conforme o padrão de \citet{bernhardt2012}.

\subsection{Fatores}

Quatro fatores binários, descritos na Tabela~\ref{tab:fatores}, alternam entre uma variante que segue o princípio (1) e outra que o viola (0). Cada variante 0 reproduz um defeito conhecido: o uso do relógio do sistema~\citep{airflow2023}, a anexação de linhas em vez da sobrescrita~\citep{airflow2023}, a mutação de dados de entrada e o argumento padrão mutável em Python, avaliado uma única vez~\citep{pythontutorial}, e a ausência de deduplicação. As 16 combinações são rotuladas pela sequência T\,W\,P\,D; por exemplo, \texttt{1101} viola apenas P.

\begin{table}[t]
  \caption{Fatores do experimento e sua classificação}
  \label{tab:fatores}
  \small
  \begin{tabular}{>{\raggedright\arraybackslash}p{0.24\columnwidth}>{\raggedright\arraybackslash}p{0.3\columnwidth}>{\raggedright\arraybackslash}p{0.33\columnwidth}}
    \toprule
    Fator & Variante 1 & Variante 0 \\
    \midrule
    T Tempo injetado \newline (funcional) & data recebida do orquestrador & data derivada do relógio (``ontem'') \\
    P Transformações puras \newline (funcional) & etapas devolvem dados novos & normalização altera o registro lido; acumulador persiste entre chamadas \\
    W Sobrescrita \newline (engenharia) & partição substituída & linhas anexadas \\
    D Deduplicação \newline (engenharia) & um registro por \texttt{id\_evento} & sem deduplicação \\
    \bottomrule
  \end{tabular}
\end{table}

Um fator é classificado como funcional quando sua variante 1 decorre de transparência referencial, isto é, quando torna a saída da etapa dependente apenas de argumentos explícitos. T torna explícita uma entrada antes oculta (o relógio), e P elimina efeitos colaterais e estado compartilhado. W e D tratam da interação com o armazenamento e da semântica dos dados. A classificação é discutível, pois \citet{beauchemin2018} inclui a sobrescrita de partições no próprio conceito de engenharia de dados funcional; ela é tratada aqui como decisão de projeto e discutida na Seção~\ref{sec:discussao}.

\subsection{Cenários, Oráculo e Métrica}

Os efeitos de P dependem de o estado em memória sobreviver entre execuções. Por isso, cada cenário declara se as execuções ocorrem em \emph{processo novo}, como uma tarefa reiniciada pelo orquestrador, ou no \emph{mesmo processo}, como uma célula reexecutada num \emph{notebook} ou um laço de \emph{backfill}. A origem mantém em cache, por processo, os registros já lidos. A Tabela~\ref{tab:cenarios} descreve os seis cenários; a execução de \emph{D} é agendada para \emph{D}+1 às 02:00.

\begin{table}[t]
  \caption{Cenários de execução}
  \label{tab:cenarios}
  \small
  \begin{tabular}{l>{\raggedright\arraybackslash}p{0.62\columnwidth}l}
    \toprule
    & Descrição & Processo \\
    \midrule
    C0 & Controle: cada dia executado uma vez & novo \\
    C1 & Execução duplicada: cada dia às 02:00 e às 02:05 & mesmo \\
    C2 & \emph{Backfill}: os 5 dias executados em sequência dias após o fim do período & mesmo \\
    C3 & Falha parcial: num dia sorteado, a escrita falha após \emph{k} linhas e a tarefa é repetida & novo \\
    C4 & Entrega duplicada: 10\% dos eventos reentregues & novo \\
    C5 & Dados atrasados: 10\% dos eventos chegam em \emph{D}+1; todos os dias são reprocessados depois & novo \\
    \bottomrule
  \end{tabular}
\end{table}

Os dados são sintéticos: 5 dias com 150 eventos cada, 40 clientes, defeitos em 15\% dos registros (campos ausentes, valores ilegíveis ou negativos) e 30 sementes. O gabarito é calculado diretamente a partir dos eventos gerados, por código que não compartilha nenhuma função com o pipeline. Ao fim de cada cenário, o conteúdo completo do destino é serializado de forma canônica e comparado com o gabarito; partições em datas erradas ou linhas duplicadas tornam o resultado incorreto. A métrica é a proporção das 30 sementes em que o destino coincide com o gabarito. C0 é condição de validade: todas as configurações devem passar nele.

\subsection{Hipóteses}

Antes das medições, previu-se analiticamente em quais cenários cada fator, violado isoladamente, falharia: T em C2 e C5 (a execução tardia processa a data errada); W em C1, C3 e C5 (a anexação duplica linhas); P em C1 e C2 (registros em cache convertidos duas vezes e acumulador com totais de execuções anteriores); e D em C4. Para as combinações, a hipótese é aditiva: uma configuração falha num cenário se algum de seus fatores violados falharia nele isoladamente.

\section{Implementação}

O núcleo usa os conceitos de programação funcional exigidos pelo estudo. A composição das etapas é definida recursivamente, com casamento de padrões sobre a tupla de etapas (Listagem~\ref{lst:compor}), e a normalização usa casamento de padrões estrutural para reconhecer os dois formatos de registro.

\begin{lstlisting}[float=t,caption={Composição recursiva das etapas},label={lst:compor}]
def compor(*etapas):
    match etapas:
        case ():
            return identidade
        case (primeira, *resto):
            seguinte = compor(*resto)
            return lambda d: seguinte(primeira(d))
\end{lstlisting}

O pipeline é montado como composição de funções de primeira classe, nomeadas e anônimas, escolhidas conforme a configuração (Listagem~\ref{lst:montar}); a agregação pura é uma dobra (\texttt{reduce}) sobre um dicionário imutável, e o processamento de coleções usa \texttt{map}, \texttt{filter}, \texttt{sorted} e \texttt{groupby}.

\begin{lstlisting}[float=t,caption={Montagem do núcleo conforme os fatores P e D (trecho adaptado)},label={lst:montar}]
def montar_transformacao(cfg, referencia, data_alvo):
    puro = cfg.transformacoes_puras
    norm = normalizar if puro else normalizar_no_lugar
    return compor(
        lambda brutos: [r for r in map(norm, brutos)
                        if r is not None],
        lambda regs: list(filter(
            lambda r: e_valido(r) and r["data"] == data_alvo,
            regs)),
        deduplicar if cfg.dedup_por_chave else identidade,
        partial(enriquecer if puro else enriquecer_no_lugar,
                referencia),
        agregar if puro
        else partial(agregar_com_estado, data_alvo=data_alvo),
    )
\end{lstlisting}

As variantes que violam P são duas: \texttt{normalizar\_no\_lugar} converte o valor para centavos no próprio registro recebido, de modo que uma segunda leitura do mesmo objeto o converte outra vez; e \texttt{agregar\_com\_estado} usa um dicionário como argumento padrão, criado uma única vez e compartilhado entre chamadas. A casca simula o início de um novo processo recarregando o módulo do núcleo e descartando o cache da origem.

\section{Resultados}

\subsection{QP1: Reexecução}

Foram realizadas 2.880 execuções (16 configurações, 6 cenários, 30 sementes). Em todas as 96 células, a proporção observada foi 0 ou 1 e coincidiu com a previsão (Tabela~\ref{tab:resultados}). C0 passou em todas as configurações, o que indica que as falhas observadas decorrem das reexecuções, e não de erros de implementação.

\begin{table}[t]
  \caption{Proporção de sementes com resultado correto (\ok{} = 1,00; \falha{} = 0,00)}
  \label{tab:resultados}
  \small
  \begin{tabular}{lcccccc}
    \toprule
    T\,W\,P\,D & C0 & C1 & C2 & C3 & C4 & C5 \\
    \midrule
    \texttt{1111} & \ok & \ok & \ok & \ok & \ok & \ok \\
    \texttt{1110} & \ok & \ok & \ok & \ok & \falha & \ok \\
    \texttt{1101} & \ok & \falha & \falha & \ok & \ok & \ok \\
    \texttt{1100} & \ok & \falha & \falha & \ok & \falha & \ok \\
    \texttt{1011} & \ok & \falha & \ok & \falha & \ok & \falha \\
    \texttt{1010} & \ok & \falha & \ok & \falha & \falha & \falha \\
    \texttt{1001} & \ok & \falha & \falha & \falha & \ok & \falha \\
    \texttt{1000} & \ok & \falha & \falha & \falha & \falha & \falha \\
    \texttt{0111} & \ok & \ok & \falha & \ok & \ok & \falha \\
    \texttt{0110} & \ok & \ok & \falha & \ok & \falha & \falha \\
    \texttt{0101} & \ok & \falha & \falha & \ok & \ok & \falha \\
    \texttt{0100} & \ok & \falha & \falha & \ok & \falha & \falha \\
    \texttt{0011} & \ok & \falha & \falha & \falha & \ok & \falha \\
    \texttt{0010} & \ok & \falha & \falha & \falha & \falha & \falha \\
    \texttt{0001} & \ok & \falha & \falha & \falha & \ok & \falha \\
    \texttt{0000} & \ok & \falha & \falha & \falha & \falha & \falha \\
    \bottomrule
  \end{tabular}
\end{table}

A tabela é explicada por um conjunto mínimo de práticas por cenário: C1 exige W e P; C2 exige T e P; C3 exige W; C4 exige D; e C5 exige T e W. As práticas funcionais (T, P) são necessárias em três dos cinco cenários de reexecução (C1, C2, C5), as de engenharia (W, D) em quatro (C1, C3, C4, C5), e os dois grupos são necessários ao mesmo tempo em C1 e C5. A configuração que adota apenas as práticas funcionais (\texttt{1010}) passa somente em C0 e C2; a que adota apenas as de engenharia (\texttt{0101}) passa em C0, C3 e C4.

Como medida exploratória, não registrada antes das medições, comparou-se a receita total gravada com a esperada. Sem W, a execução duplicada (C1) grava exatamente o dobro da receita. Sem T, o \emph{backfill} (C2) grava todas as execuções na partição de um único dia sem eventos, e as partições dos cinco dias processados nunca são gravadas. Sem P, a receita em C1 é 101 vezes a esperada; esse valor é um artefato do defeito escolhido (a reconversão multiplica o valor por 100, e o acumulador soma a primeira execução) e indica apenas que o erro é grande e silencioso, não sua magnitude típica.

\subsection{QP2: Testabilidade}

\emph{Em andamento.} O protocolo será fixado numa nova versão da especificação antes da escrita das suítes de teste. A intenção é comparar as variantes pura e impura do núcleo pelo número de dublês de teste necessários e pelo escore de mutação~\citep{jia2011}, com propriedades verificadas por meio da Hypothesis~\citep{maciver2019}.

\section{Discussão e Ameaças à Validade}
\label{sec:discussao}

\emph{Versão preliminar.} O resultado principal da QP1 é que a recomendação ``escreva o pipeline em estilo funcional'' é incompleta: nenhum dos dois grupos de práticas garante a reexecução segura sozinho. As práticas classificadas como funcionais não protegem contra entrega duplicada nem contra falhas parciais de escrita, e as de engenharia não protegem contra a data derivada do relógio nem contra estado compartilhado. A literatura normativa já recomenda as quatro práticas em conjunto~\citep{beauchemin2018,airflow2023}; o que este estudo acrescenta é o mapeamento verificável entre cada prática e as falhas que ela evita.

Quatro ameaças limitam essa conclusão. \emph{Construção:} o autor escreveu todas as variantes, e a variante 0 de cada fator pode não representar o código encontrado na prática; por isso, os resultados são em grande parte previsíveis por análise, e a contribuição está na verificação executável do mapeamento e da hipótese aditiva, não na descoberta de falhas inesperadas. \emph{Interna:} a necessidade de P depende do modelo de processo; as variantes impuras deste pipeline só falharam nos cenários em que o estado sobrevive entre execuções (C1, C2), e não falhariam neles se cada execução ocorresse em processo novo. \emph{Classificação:} se a sobrescrita de partição for considerada funcional, como em \citet{beauchemin2018}, as práticas funcionais passam a ser necessárias em todos os cenários exceto C4. \emph{Externa:} trata-se de um único pipeline, pequeno, sem paralelismo e em Python puro.

\section{Conclusão}

\emph{Versão preliminar.} Num desenho fatorial com 16 configurações e seis cenários, as falhas de reexecução de um pipeline em lote seguiram exatamente a previsão aditiva, e cada cenário exigiu um conjunto específico de práticas. A pureza e o tempo injetado são necessários, mas não suficientes, para a idempotência. Os próximos passos são a avaliação da QP2 e a discussão das implicações para pipelines em motores distribuídos.

\section*{Disponibilidade de Artefatos}

O código do pipeline, o gerador de dados, o oráculo, os testes, a especificação do experimento e os resultados brutos estão disponíveis em [LINK A DEFINIR], sob a licença [A DEFINIR]. A \emph{tag} \texttt{especificacao-v1.0} marca a versão da especificação registrada antes das medições.
